\section{Gestión de requisitos y trazabilidad}
\label{sec:gestion}

\subsection{Línea base y control de versiones}

El ERS se versiona en el repositorio del equipo junto con sus fuentes \LaTeX{}, de modo que el
documento entregado se regenera por clonación y compilación. La línea base aprobada por el Change
Control Board al cierre de la Unidad IV se publicó como etiqueta anotada \id{baseline-v1.1}, y la
Unidad V cierra con la etiqueta \id{baseline-final}.

\begin{table}[H]
\centering\footnotesize
\caption{Historial de versiones del ERS}
\label{tab:versiones}
\begin{tabular}{|C{2.0cm}|C{2.2cm}|C{2.4cm}|L{7.4cm}|}
\hline
\rowcolor{grisclaro}\textbf{Versión} & \textbf{Rev. interna} & \textbf{Etiqueta Git} & \textbf{Contenido}\\ \hline
v1.0 & 3.0 & --- & Entrega 2A. Documento sometido a inspección formal\\ \hline
v1.1 & 3.1 & \id{baseline-v1.1}, \id{baseline-final} & Correcciones de la inspección, cambios aprobados por el CCB, y cierre de la Unidad V: auditoría, DFD, casos de uso, criterios BDD y componentes de IA\\ \hline
\end{tabular}
\end{table}

\noindent
La numeración merece una aclaración porque conviven dos. El ERS registra internamente sus revisiones
desde la Entrega 1A y llegó a la Entrega 2A como revisión interna 3.0; la asignatura designa a ese
mismo documento como v1.0. Se adopta la numeración de la asignatura como versión oficial, declarada
de forma idéntica en la portada del ERS, en su historial, en el \id{CHANGELOG.md} y en la etiqueta de
Git, y el historial del documento rotula su columna como \emph{revisión interna} para que la
coexistencia no produzca ambigüedad.

\subsection{Control de cambios}

Las tres solicitudes tramitadas ante el Change Control Board se recogen en la
\S\ref{sec:validacion}. El procedimiento seguido ---formulario con análisis de impacto derivado de la
matriz, deliberación con roles diferenciados, decisión motivada y acta firmada--- responde al modelo
de control integrado de cambios del PMBOK \cite{pmbok7}.

El análisis de impacto de cada solicitud se obtuvo recorriendo la matriz de trazabilidad y filtrando
por el requisito afectado, no por estimación. Esa diferencia es la que separa un análisis de impacto
de una conjetura informada: \id{RFC-03} declaró afectar a las filas \id{TR-03} y \id{TR-34} y añadir
\id{TR-53} a \id{TR-55}, y esos identificadores son verificables en el artefacto.

% =====================================================================
%  Sección de trazabilidad extremo a extremo — informe PE5
%  Requiere: longtable, array, xcolor, colortbl, float, y la macro \id
% =====================================================================

\subsection{Trazabilidad extremo a extremo}

\subsubsection*{Cadena de trazabilidad adoptada}

La trazabilidad de las entregas anteriores enlazaba la parte interesada con el requisito, el caso de
uso y el prototipo. Esa cadena bastaba para responder de dónde viene un requisito y qué pantalla lo
materializa, pero no para responder qué proceso lo ejecuta, sobre qué estado del dominio actúa ni
con qué condición se acepta. La Unidad~V extiende la cadena a ocho eslabones:

\begin{center}
\footnotesize
Fuente $\rightarrow$ Requisito $\rightarrow$ Caso de uso $\rightarrow$ Clase/Módulo $\rightarrow$
Proceso DFD $\rightarrow$ Estado $\rightarrow$ Criterio BDD $\rightarrow$ Caso de prueba
\end{center}

\noindent
Tres eslabones son nuevos y ninguno existía como columna antes de esta práctica:

\begin{itemize}[leftmargin=*, nosep]
  \item \textbf{Proceso DFD.} No existía porque el ERS no tenía diagramas de flujo de datos. Se
        incorporaron en la \S8 del ERS con siete procesos, uno por bloque funcional, de modo que la
        columna se deriva del bloque en que el propio documento ya ubicaba cada requisito y no
        introduce una clasificación nueva.
  \item \textbf{Estado.} Enlaza el requisito con la transición de la máquina de estados sobre la que
        actúa. El ERS modela el ciclo de vida de dos entidades: el racimo y la alerta.
  \item \textbf{Criterio BDD.} Enlaza el requisito con su condición de aceptación en formato
        Dado--Cuando--Entonces: \id{CA-01} a \id{CA-24} cuando procede de una historia de usuario y
        \id{CB-01} a \id{CB-37} cuando se deriva de la ficha del requisito.
\end{itemize}

\subsubsection*{Criterio de completitud de una traza}

Una fila se declara \textbf{completa} cuando tiene los seis eslabones exigibles: fuente, caso de uso,
módulo, proceso, criterio de aceptación y caso de prueba. Se declara \textbf{parcial} si falta uno y
\textbf{huérfana} si faltan dos o más.

El eslabón de estado no entra en ese conteo, y la razón conviene declararla porque afecta a la
lectura de la métrica. Solo dos entidades del dominio tienen ciclo de vida modelado; los requisitos
cuyo objeto no es un racimo ni una alerta ---gestión de personal, planificación de labores,
generación de reportes--- no tienen estado asociado. Exigirles uno obligaría a inventar máquinas de
estados sin sustento en el dominio para que una métrica saliera mejor. La columna declara
\emph{No aplica} en esas 34 filas, y esa declaración es información, no un vacío.

\subsubsection*{Matriz completa}

{\scriptsize
\begin{longtable}{|M{0.7cm}|L{1.9cm}|M{1.3cm}|M{1.0cm}|M{1.0cm}|M{1.1cm}|L{3.2cm}|M{1.1cm}|M{1.0cm}|M{1.3cm}|}
\caption{Matriz de trazabilidad extremo a extremo (73 filas). El asterisco indica que el proceso se deriva del caso de uso y no del bloque funcional, por tratarse de un requisito no funcional o de una restricción. En la columna de estado, \textbf{R} designa la máquina de estados del racimo y \textbf{A} la de la alerta.}
\label{tab:matriz-e2e}\\
\hline
\rowcolor{grisclaro}\textbf{N.º} & \textbf{Fuente} & \textbf{Requisito} & \textbf{CU} & \textbf{Módulo} & \textbf{Proceso} & \textbf{Estado} & \textbf{BDD} & \textbf{Prueba} & \textbf{Traza}\\ \hline
\endfirsthead
\hline
\rowcolor{grisclaro}\textbf{N.º} & \textbf{Fuente} & \textbf{Requisito} & \textbf{CU} & \textbf{Módulo} & \textbf{Proceso} & \textbf{Estado} & \textbf{BDD} & \textbf{Prueba} & \textbf{Traza}\\ \hline
\endhead
1 & Administrador & \id{RF-01} & CU-11 & M-01 & P-01 & No aplica & CA-01 & CP-01 & Completa\\ \hline
2 & Administrador & \id{RF-02} & CU-12 & M-01 & P-01 & No aplica & CA-02 & CP-02 & Completa\\ \hline
3 & Administrador & \id{RF-03} & CU-13 & M-01 & P-01 & No aplica & CA-03 & CP-03 & Completa\\ \hline
4 & Trabajador & \id{RF-04} & CU-04 & M-03 & P-03 & No aplica & CA-04 & CP-04 & Completa\\ \hline
5 & Trabajador & \id{RF-04} & CU-04 & M-03 & P-03 & No aplica & CA-04 & CP-04 & Completa\\ \hline
6 & Técnico fitosan. & \id{RF-05} & CU-15 & M-03 & P-03 & A: Abierta $\to$ En verificación & CA-05 & CP-05 & Completa\\ \hline
7 & Administrador & \id{RF-06} & CU-12 & M-03 & P-03 & R: En desarrollo & CB-01 & CP-06 & Completa\\ \hline
8 & Asistente admin. & \id{RF-07} & CU-01 & M-05 & P-05 & A: Abierta $\to$ Notificada & CA-06 & CP-07 & Completa\\ \hline
9 & Asistente admin. & \id{RF-07} & CU-01 & M-05 & P-05 & A: Abierta $\to$ Notificada & CA-06 & CP-07 & Completa\\ \hline
10 & Asesor técnico & \id{RF-08} & CU-02 & M-05 & P-05 & A: Abierta $\to$ Notificada & CA-07 & CP-08 & Completa\\ \hline
11 & Asesor técnico & \id{RF-08} & CU-02 & M-05 & P-05 & A: Abierta $\to$ Notificada & CA-07 & CP-08 & Completa\\ \hline
12 & Técnico fitosan. & \id{RF-09} & CU-01 & M-05 & P-05 & A: En verificación $\to$ Atendida & CB-02 & CP-09 & Completa\\ \hline
13 & Asesor técnico & \id{RF-10} & CU-12 & M-01 & P-01 & No aplica & CA-08 & CP-10 & Completa\\ \hline
14 & Asesor técnico & \id{RF-11} & CU-02 & M-03 & P-03 & No aplica & CB-03 & CP-11 & Completa\\ \hline
15 & Administrador & \id{RF-12} & CU-16 & M-05 & P-05 & A: (inicial) $\to$ Abierta & CA-09 & CP-12 & Completa\\ \hline
16 & Trabajador & \id{RF-12} & CU-16 & M-05 & P-05 & A: (inicial) $\to$ Abierta & CA-09 & CP-12 & Completa\\ \hline
17 & Jefe polinización & \id{RF-13} & CU-03 & M-04 & P-04 & R: Inflorescencia $\to$ Polinizado & CA-10 & CP-13 & Completa\\ \hline
18 & Jefe polinización & \id{RF-14} & CU-03 & M-04 & P-04 & R: Inflorescencia $\to$ Polinizado & CA-11 & CP-14 & Completa\\ \hline
19 & Administrador & \id{RF-15} & CU-03 & M-04 & P-04 & No aplica & CB-04 & CP-15 & Completa\\ \hline
20 & Administrador & \id{RF-16} & CU-05 & M-04 & P-04 & No aplica & CB-05 & CP-16 & Completa\\ \hline
21 & Administrador & \id{RF-17} & CU-05 & M-03 & P-03 & No aplica & CB-06 & CP-17 & Completa\\ \hline
22 & Administrador & \id{RF-18} & CU-17 & M-07 & P-07 & R: En desarrollo & CA-12 & CP-18 & Completa\\ \hline
23 & Administrador & \id{RF-19} & CU-05 & M-07 & P-07 & No aplica & CA-13 & CP-19 & Completa\\ \hline
24 & Administrador & \id{RF-20} & CU-05 & M-07 & P-07 & No aplica & CB-07 & CP-20 & Completa\\ \hline
25 & Téc. extractora & \id{RF-21} & CU-06 & M-05 & P-05 & R: En desarrollo $\to$ Maduro & CA-14 & CP-21 & Completa\\ \hline
26 & Téc. extractora & \id{RF-21} & CU-06 & M-05 & P-05 & R: En desarrollo $\to$ Maduro & CA-14 & CP-21 & Completa\\ \hline
27 & Administrador & \id{RF-22} & CU-06 & M-06 & P-06 & R: Maduro $\to$ Cosechado & CA-15 & CP-22 & Completa\\ \hline
28 & Administrador & \id{RF-23} & CU-10 & M-06 & P-06 & R: Despachado $\to$ Calificado & CB-08 & CP-23 & Completa\\ \hline
29 & Téc. extractora & \id{RF-24} & CU-10 & M-06 & P-06 & R: Calificado & CB-09 & CP-24 & Completa\\ \hline
30 & Téc. extractora & \id{RF-25} & CU-10 & M-06 & P-06 & R: Cosechado $\to$ Degradado & CB-10 & CP-25 & Completa\\ \hline
31 & Administrador & \id{RF-26} & CU-07 & M-02 & P-02 & No aplica & CA-16 & CP-26 & Completa\\ \hline
32 & Administrador & \id{RF-27} & CU-07 & M-02 & P-02 & No aplica & CB-11 & CP-27 & Completa\\ \hline
33 & Trabajador & \id{RF-28} & CU-04 & M-03 & P-03 & No aplica & CA-17 & CP-28 & Completa\\ \hline
34 & Trabajador & \id{RF-29} & CU-08 & M-03 & P-03 & No aplica & CB-12 & CP-29 & Completa\\ \hline
35 & Asistente admin. & \id{RF-30} & CU-09 & M-03 & P-03 & A: Abierta $\to$ En verificación & CA-18 & CP-30 & Completa\\ \hline
36 & Asistente admin. & \id{RF-31} & CU-09 & M-03 & P-03 & A: En verificación $\to$ Atendida & CB-13 & CP-31 & Completa\\ \hline
37 & Téc. extractora & \id{RF-32} & CU-17 & M-07 & P-07 & No aplica & CB-14 & CP-32 & Completa\\ \hline
38 & Téc. extractora & \id{RF-33} & CU-10 & M-06 & P-06 & R: Cosechado $\to$ Despachado & CB-15 & CP-33 & Completa\\ \hline
39 & Administrador & \id{RF-34} & CU-07 & M-02 & P-02 & No aplica & CB-16 & CP-34 & Completa\\ \hline
40 & Jefe de campo & \id{RF-35} & CU-04 & M-02 & P-02 & No aplica & CA-19 & CP-35 & Completa\\ \hline
41 & Téc. extractora & \id{RNF-01} & CU-06 & M-05 & P-05$^{*}$ & No aplica & CB-19 & CP-36 & Completa\\ \hline
42 & Trabajador & \id{RNF-11} & CU-04 & M-03 & P-03$^{*}$ & No aplica & CB-29 & CP-37 & Completa\\ \hline
43 & Jefe polinización & \id{RNF-05} & CU-03 & M-04 & P-04$^{*}$ & No aplica & CB-23 & CP-38 & Completa\\ \hline
44 & Trabajador & \id{RNF-08} & CU-04 & M-03 & P-03$^{*}$ & No aplica & CB-26 & CP-39 & Completa\\ \hline
45 & Trabajador & \id{RNF-14} & CU-03 & M-04 & P-04$^{*}$ & No aplica & CB-32 & CP-40 & Completa\\ \hline
46 & Asistente admin. & \id{RNF-16} & CU-18 & M-05 & P-05$^{*}$ & No aplica & CB-34 & CP-41 & Completa\\ \hline
47 & Trabajador & \id{RD-10} & CU-04 & M-03 & P-03$^{*}$ & No aplica & No aplica (restricción) & CP-42 & Completa\\ \hline
48 & Téc. extractora & \id{RD-07} & CU-06 & M-05 & P-05$^{*}$ & No aplica & No aplica (restricción) & CP-43 & Completa\\ \hline
49 & Administrador & \id{RF-36} & CU-08 & M-02 & P-02 & No aplica & CA-20 & CP-44 & Completa\\ \hline
50 & Administrador & \id{RF-37} & CU-07 & M-02 & P-02 & No aplica & CA-21 & CP-45 & Completa\\ \hline
51 & Administrador & \id{RF-38} & CU-07 & M-02 & P-02 & No aplica & CB-17 & CP-46 & Completa\\ \hline
52 & Administrador & \id{RF-39} & CU-07 & M-02 & P-02 & No aplica & CB-18 & CP-47 & Completa\\ \hline
53 & Trabajador & \id{RF-40} & CU-13 & M-03 & P-03 & No aplica & CA-22 & CP-53 & Completa\\ \hline
54 & Trabajador & \id{RF-41} & CU-13 & M-03 & P-03 & No aplica & CA-23 & CP-54 & Completa\\ \hline
55 & Administrador & \id{RF-42} & CU-13 & M-03 & P-03 & No aplica & CA-24 & CP-55 & Completa\\ \hline
56 & Asist. admin. & \id{RF-IA-01} & CU-01 & M-05 & P-05 & A: Abierta $\to$ Notificada & en ficha & CP-56 & Completa\\ \hline
57 & Asist. admin. & \id{RF-IA-02} & CU-01 & M-05 & P-05 & No aplica & en ficha & CP-57 & Completa\\ \hline
58 & Admin. & \id{RF-IA-03} & CU-18 & M-05 & P-05 & A: $\to$ Descartada & en ficha & CP-58 & Completa\\ \hline
59 & Téc. extr. & \id{RF-IA-04} & CU-06 & M-06 & P-06 & R: En desarrollo $\to$ Maduro & en ficha & CP-59 & Completa\\ \hline
60 & Téc. extr. & \id{RF-IA-05} & CU-06 & M-06 & P-06 & R: Maduro $\to$ Cosechado & en ficha & CP-60 & Completa\\ \hline
61 & Admin. & \id{RF-IA-06} & CU-10 & M-06 & P-06 & R: Cosechado $\to$ Despachado & en ficha & CP-61 & Completa\\ \hline
62 & Asist. admin. & \id{RNF-IA-01} & CU-01 & M-05 & P-05$^{*}$ & No aplica & en ficha & CP-62 & Completa\\ \hline
63 & Asist. admin. & \id{RNF-IA-02} & CU-01 & M-05 & P-05$^{*}$ & No aplica & en ficha & CP-63 & Completa\\ \hline
64 & Trabajador & \id{RNF-IA-03} & CU-01 & M-05 & P-05$^{*}$ & No aplica & en ficha & CP-64 & Completa\\ \hline
65 & Admin. & \id{RNF-IA-04} & CU-01 & M-05 & P-05$^{*}$ & No aplica & en ficha & CP-65 & Completa\\ \hline
66 & Trabajador & \id{RNF-IA-05} & CU-01 & M-05 & P-05$^{*}$ & No aplica & en ficha & CP-66 & Completa\\ \hline
67 & Asist. admin. & \id{RNF-IA-06} & CU-18 & M-05 & P-05$^{*}$ & No aplica & en ficha & CP-67 & Completa\\ \hline
68 & Téc. extr. & \id{RNF-IA-07} & CU-06 & M-06 & P-06$^{*}$ & No aplica & en ficha & CP-68 & Completa\\ \hline
69 & Téc. extr. & \id{RNF-IA-08} & CU-06 & M-06 & P-06$^{*}$ & No aplica & en ficha & CP-69 & Completa\\ \hline
70 & Admin. & \id{RNF-IA-09} & CU-06 & M-06 & P-06$^{*}$ & No aplica & en ficha & CP-70 & Completa\\ \hline
71 & Trabajador & \id{RNF-IA-10} & CU-10 & M-06 & P-06$^{*}$ & No aplica & en ficha & CP-71 & Completa\\ \hline
72 & Admin. & \id{RNF-IA-11} & CU-10 & M-06 & P-06$^{*}$ & No aplica & en ficha & CP-72 & Completa\\ \hline
73 & Admin. & \id{RNF-IA-12} & CU-18 & M-06 & P-06$^{*}$ & No aplica & en ficha & CP-73 & Completa\\ \hline
\end{longtable}}

\subsubsection*{Diagnóstico y elementos huérfanos}

\begin{table}[H]
\centering\footnotesize
\caption{Diagnóstico de la trazabilidad}
\label{tab:diag-traza}
\begin{tabular}{|L{7.0cm}|C{2.4cm}|L{6.2cm}|}
\hline
\rowcolor{grisclaro}\textbf{Indicador} & \textbf{Valor} & \textbf{Lectura}\\ \hline
Filas de la matriz & 73 & Cubre 42 RF, 10 RNF, 3 RD y los 18 requisitos de IA\\ \hline
Trazas completas & 73 (100\,\%) & Ninguna fila carece de un eslabón exigible\\ \hline
Trazas parciales & 0 & ---\\ \hline
Trazas huérfanas & 0 & ---\\ \hline
Filas con estado modelado & 26 & Las 47 restantes declaran \emph{No aplica}\\ \hline
Filas con criterio BDD & 71 & Las 2 restantes son restricciones de diseño\\ \hline
Filas sin prototipo & 1 & Requisito de prioridad distinta de \emph{Must}\\ \hline
\end{tabular}
\end{table}

\noindent
La matriz no arroja elementos huérfanos, pero sí tres vacíos que se declaran de forma expresa en
lugar de rellenarlos:

\begin{enumerate}[leftmargin=*, nosep]
  \item \textbf{Requisitos sin prototipo de interfaz.} \id{RF-03}, \id{RF-27} y \id{RF-34}. Los tres
        son de prioridad distinta de \emph{Must}, de modo que el vacío no incumple la exigencia de
        prototipar las pantallas críticas. Se difiere a la Entrega~4.
  \item \textbf{Restricciones sin criterio de aceptación.} \id{RD-07} y \id{RD-10}. Una restricción
        de diseño no describe un comportamiento del sistema sino una condición impuesta sobre su
        construcción, y por tanto no admite criterio de aceptación en el mismo sentido que un
        requisito. Quedan fuera del denominador de la métrica de verificabilidad.
  \item \textbf{Requisitos de IA sin criterio BDD en columna propia.} Los dieciocho requisitos de
        inteligencia artificial llevan su criterio de aceptación dentro de la propia ficha de la \S9
        del ERS, en formato Dado--Cuando--Entonces, y por eso la columna remite a ella en lugar de
        citar un identificador \id{CA-XX} o \id{CB-XX}. La decisión evita duplicar el criterio en dos
        lugares del documento, que es precisamente el mecanismo que produjo los defectos de
        consistencia detectados en la Unidad~IV.
  \item \textbf{Requisitos no funcionales sin proceso propio.} Un requisito no funcional no realiza
        un proceso: lo condiciona. Se les asigna el proceso mayoritario de los requisitos funcionales
        que comparten su caso de uso, y la matriz lo marca con asterisco para que la derivación sea
        visible y no se confunda con una asignación directa.
\end{enumerate}

\subsubsection*{Sincronización entre el ERS y el tablero CASE}

La trazabilidad solo es útil si el tablero refleja el documento. El contraste cuantitativo entre
ambos artefactos, antes y después de la actualización del backlog, es el siguiente:

\begin{table}[H]
\centering\footnotesize
\caption{Sincronización entre el ERS y el tablero de Jira}
\label{tab:sincro}
\begin{tabular}{|L{4.6cm}|C{1.5cm}|C{1.6cm}|C{1.6cm}|C{1.6cm}|C{1.6cm}|}
\hline
\rowcolor{grisclaro}\textbf{Elemento} & \textbf{En el ERS} & \textbf{Tablero inicial} & \textbf{Diverg. inicial} & \textbf{Tablero final} & \textbf{Diverg. final}\\ \hline
Epics / bloques funcionales & 10 & 7 & 3 & 10 & 0\\ \hline
Stories / requisitos & 79 & 39 & 40 & 79 & 0\\ \hline
Sub-tasks / criterios de aceptación & 79 & 38 & 41 & 86 & 0\\ \hline
Historias de usuario & 24 & 19 & 5 & 24 & 0\\ \hline
Filas de trazabilidad & 73 & 52 & 21 & 73 & 0\\ \hline
\multicolumn{3}{|r|}{\textbf{Índice de sincronización}} & \textbf{58,5\,\%} & & \textbf{100\,\%}\\ \hline
\end{tabular}
\end{table}

\noindent
La cifra publicada al cierre de la Unidad~IV fue 67,1\,\%. Se conserva aquí porque es la que consta
en aquel informe, pero el denominador con que se calculó era incompleto: contabilizaba 48 elementos
especificados donde el ERS especifica 79, porque no computaba los diecinueve requisitos no
funcionales ni los doce \id{RNF-IA} como elementos que el tablero debía reflejar. Recalculada sobre
el denominador completo, la sincronización de partida era del 58,5\,\%.

Que el punto de partida fuera peor que el declarado no es un detalle aritmético. Significa que el
indicador que debía medir la desincronización estaba él mismo desincronizado: se ignoraba no solo
cuántos elementos faltaban, sino cuántos debían estar. Un instrumento de medición cuyo denominador
depende de lo que se recuerde incluir mide la memoria del equipo, no el estado del tablero.

La divergencia inicial tiene dos causas verificables. La primera, que la carga del backlog se
ejecutó en la Unidad~IV, antes de que el Change Control Board aprobara \id{RFC-01} a \id{RFC-03}:
los cambios se implementaron en el documento y no se propagaron al tablero. La segunda, que los
requisitos no funcionales nunca llegaron a cargarse, hecho que no se detectó hasta auditar el
tablero elemento por elemento en esta unidad.

El hallazgo es más interesante que el número. Un backlog desincronizado no es un descuido puntual:
es la consecuencia previsible de mantener dos representaciones del mismo conjunto de requisitos sin
un mecanismo que las reconcilie. La misma causa produjo, en la Unidad~IV, once de los treinta y tres
defectos de consistencia detectados en la inspección. La acción correctiva registrada no es solo
cargar los elementos que faltan, sino incorporar al procedimiento de cierre de cada RFC la
actualización del tablero como paso obligatorio, del mismo modo que ya lo es la actualización del
\id{CHANGELOG.md}.

Al cierre de esta unidad la sincronización es del 100\,\%: los 175 elementos del tablero ---diez
epics, setenta y nueve \emph{stories} y ochenta y seis \emph{sub-tasks}--- cubren la totalidad de lo
especificado. Las siete \emph{sub-tasks} que exceden a los criterios del documento son tareas de
verificación creadas por el equipo para su propio seguimiento y no derivan de un criterio de
aceptación, por lo que no se computan como divergencia: el índice mide si algo especificado quedó
sin cargar, no si el tablero contiene elementos adicionales. El respaldo es la exportación
\id{backlog\_export.csv} del Apéndice de trazabilidad.
